\documentclass[11pt,letterpaper]{article}

\usepackage[letterpaper,top=0.5in,bottom=0.5in,left=0.55in,right=0.55in]{geometry}
\usepackage{graphicx}
\usepackage{xcolor}
\usepackage{listings}
\usepackage{booktabs}
\usepackage{hyperref}
\usepackage{caption}
\usepackage{float}
\usepackage{titlesec}

\hypersetup{
    colorlinks=true,
    linkcolor=blue!60!black,
    urlcolor=blue!60!black,
}

\setlength{\parskip}{2pt}
\setlength{\parindent}{0pt}
\titlespacing*{\section}{0pt}{6pt}{2pt}
\titlespacing*{\subsection}{0pt}{4pt}{2pt}
\captionsetup{skip=2pt,font=small}

% -----------------------------------------------------------------------------
% Code listing style (C)
% -----------------------------------------------------------------------------
\definecolor{codebg}{RGB}{248,248,248}
\definecolor{codekw}{RGB}{0,0,180}
\definecolor{codecmt}{RGB}{0,120,0}
\definecolor{codestr}{RGB}{170,55,20}
\definecolor{codenum}{RGB}{120,120,120}

\lstset{
    language=C,
    basicstyle=\ttfamily\footnotesize,
    keywordstyle=\color{codekw}\bfseries,
    commentstyle=\color{codecmt}\itshape,
    stringstyle=\color{codestr},
    numbers=left,
    numberstyle=\tiny\color{codenum},
    stepnumber=1,
    numbersep=5pt,
    backgroundcolor=\color{codebg},
    frame=single,
    rulecolor=\color{black!30},
    tabsize=4,
    showstringspaces=false,
    breaklines=true,
    breakatwhitespace=false,
    breakindent=0pt,
    postbreak=\mbox{\textcolor{red}{$\hookrightarrow$}\space},
    captionpos=b,
    aboveskip=4pt,
    belowskip=4pt,
    xleftmargin=6pt,
    xrightmargin=2pt,
    columns=fullflexible,
    keepspaces=true,
    morekeywords={u32,u64,XGpio,seq_state_t}
}

\title{\vspace{-0.6in}\textbf{ECEC 661 -- Midterm Submission}\\
       \large Button Sequence Detector (BTN0, BTN0, BTN1) on Cora Z7-07S}
\author{Gary Pham \texttt{(gp492)}}
\date{\today}

\begin{document}
\maketitle
\vspace{-0.25in}

% =============================================================================
\section{Problem Statement}
% =============================================================================

Build a Zynq PS+PL system on the Digilent Cora Z7-07S that detects the
button-press sequence \textbf{BTN0, BTN0, BTN1}. A bare-metal C
application running on the PS reads the buttons and lights an on-board
LED to indicate detection. Two AXI GPIO slaves bridge the PS to the
buttons (input) and to the RGB LEDs (output); no custom IP is required.

% =============================================================================
\section{Block Design}
% =============================================================================

The block design (Figure~\ref{fig:bd}) instantiates the
\texttt{ZYNQ7 Processing System} (UART0 on MIO 14/15 for the
FT2232HQ USB-UART), a \texttt{Processor System Reset}, an
\texttt{AXI Interconnect} (1 slave, 2 masters), and two
\texttt{AXI GPIO} slaves: \texttt{axi\_gpio\_0} drives the 6-bit RGB
LED vector \texttt{rgb\_leds\_tri\_o[5:0]}, \texttt{axi\_gpio\_1}
samples the 2-bit button vector \texttt{btns\_2bits\_tri\_i[1:0]}.
The PS clock (\texttt{FCLK\_CLK0}, 100\,MHz) and reset
(\texttt{peripheral\_aresetn}) feed every AXI port. Address map:
\texttt{axi\_gpio\_0 = 0x4120\_0000}, \texttt{axi\_gpio\_1 = 0x4121\_0000}.

\begin{figure}[H]
    \centering
    \includegraphics[width=\textwidth]{diagram.png}
    \caption{Vivado block diagram \texttt{seq\_det.bd}: Zynq PS, AXI
             Interconnect, and the two AXI GPIO peripherals exposed
             externally as \texttt{rgb\_leds} and \texttt{btns\_2bits}.}
    \label{fig:bd}
\end{figure}

% =============================================================================
\section{Detection FSM}
% =============================================================================

The detector is a 4-state Moore-style machine implemented in software.
The detection LED (\texttt{LD0\_G} = \texttt{rgb\_leds\_tri\_o[1]}) is
driven directly from the current state -- it is ON if and only if
\texttt{state == S3}. KMP-style transitions ensure that a redundant
\texttt{BTN0} does not throw away the prefix already matched.

\begin{table}[H]
    \centering
    \caption{Transition table for pattern \texttt{BTN0, BTN0, BTN1}.}
    \label{tab:fsm}
    \small
    \begin{tabular}{cccl}
        \toprule
        \textbf{State} & \textbf{Event} & \textbf{Next} & \textbf{Note} \\
        \midrule
        S0 & BTN0 & S1 & start of pattern                                  \\
        S0 & BTN1 & S0 & ignore                                             \\
        S1 & BTN0 & S2 & two-in-a-row                                       \\
        S1 & BTN1 & S0 & reset                                              \\
        S2 & BTN0 & S2 & KMP self-loop (extra BTN0 is a new candidate)      \\
        S2 & BTN1 & S3 & \textbf{detected} -- LED ON                        \\
        S3 & BTN0 & S1 & leave detect; this BTN0 may open the next match    \\
        S3 & BTN1 & S0 & leave detect; reset                                \\
        \bottomrule
    \end{tabular}
\end{table}

% =============================================================================
\section{Software (\texttt{main.c})}
% =============================================================================

The PS application uses the \texttt{XGpio} driver, polls the buttons
every 5\,ms with rising-edge detection, advances the FSM, drives the
LED from the resulting state, and prints every transition over UART
for grading evidence.

\begin{lstlisting}[caption={\texttt{Midterm/real/sw/main.c}},label={lst:main}]
#include "xparameters.h"
#include "xgpio.h"
#include "xil_printf.h"
#include "xil_types.h"
#include "sleep.h"

#ifdef XPAR_AXI_GPIO_0_BASEADDR
#define LEDS_BASEADDR   XPAR_AXI_GPIO_0_BASEADDR
#else
#define LEDS_BASEADDR   0x41200000U
#endif
#ifdef XPAR_AXI_GPIO_1_BASEADDR
#define BTNS_BASEADDR   XPAR_AXI_GPIO_1_BASEADDR
#else
#define BTNS_BASEADDR   0x41210000U
#endif

#define GPIO_CHANNEL    1U
#define LED_MASK        0x3FU
#define LED_DETECT_BIT  (1U << 1)   /* LD0_G */
#define BTN0_MASK       (1U << 0)
#define BTN1_MASK       (1U << 1)
#define POLL_PERIOD_US  5000U

typedef enum { S0 = 0, S1 = 1, S2 = 2, S3 = 3 } seq_state_t;

static const char *state_name(seq_state_t s) {
    switch (s) { case S0: return "S0"; case S1: return "S1";
                 case S2: return "S2"; case S3: return "S3";
                 default: return "??"; }
}

/* 4-state KMP transition table (S3 = detected, LED ON) */
static seq_state_t fsm_step(seq_state_t s, int btn_idx) {
    if (btn_idx == 0) {
        switch (s) { case S0: return S1; case S1: return S2;
                     case S2: return S2; case S3: return S1; }
    } else {
        switch (s) { case S0: return S0; case S1: return S0;
                     case S2: return S3; case S3: return S0; }
    }
    return S0;
}

int main(void) {
    XGpio led_gpio, btn_gpio;
    if (XGpio_Initialize(&led_gpio, LEDS_BASEADDR) != XST_SUCCESS) return XST_FAILURE;
    if (XGpio_Initialize(&btn_gpio, BTNS_BASEADDR) != XST_SUCCESS) return XST_FAILURE;
    XGpio_SetDataDirection(&led_gpio, GPIO_CHANNEL, 0x00);        /* outputs */
    XGpio_SetDataDirection(&btn_gpio, GPIO_CHANNEL, 0xFFFFFFFF);  /* inputs  */
    XGpio_DiscreteWrite(&led_gpio, GPIO_CHANNEL, 0);

    for (int i = 3; i > 0; i--) {
        xil_printf("seq_det starting in %d...\r\n", i); sleep(1);
    }
    xil_printf("\r\n=== ECEC 661 Midterm: pattern BTN0,BTN0,BTN1 (4-state FSM) ===\r\n");

    seq_state_t state = S0;
    u32 prev_btns = XGpio_DiscreteRead(&btn_gpio, GPIO_CHANNEL) & 0x3U;
    unsigned detect_count = 0;

    while (1) {
        u32 cur_btns = XGpio_DiscreteRead(&btn_gpio, GPIO_CHANNEL) & 0x3U;
        if (cur_btns == 0x3U) {                 /* both held: ambiguous, skip */
            usleep(POLL_PERIOD_US); prev_btns = cur_btns; continue;
        }
        u32 rising = cur_btns & ~prev_btns;     /* 0 -> 1 edges only */
        prev_btns = cur_btns;

        int btn_idx = -1;
        if      (rising & BTN0_MASK) btn_idx = 0;
        else if (rising & BTN1_MASK) btn_idx = 1;

        if (btn_idx >= 0) {
            seq_state_t prev = state;
            state = fsm_step(state, btn_idx);

            /* Moore output: LED on iff state == S3 */
            u32 led_state = ((state == S3) ? LED_DETECT_BIT : 0U) & LED_MASK;
            XGpio_DiscreteWrite(&led_gpio, GPIO_CHANNEL, led_state);

            int just_detected = (state == S3 && prev != S3);
            if (just_detected) detect_count++;
            if (just_detected) {
                xil_printf("[BTN%d] %s -> %s  LED=ON   *** DETECTED #%u ***\r\n",
                           btn_idx, state_name(prev), state_name(state), detect_count);
            } else {
                xil_printf("[BTN%d] %s -> %s  LED=%s\r\n",
                           btn_idx, state_name(prev), state_name(state),
                           (state == S3) ? "ON " : "OFF");
            }
        }
        usleep(POLL_PERIOD_US);
    }
    return 0;
}
\end{lstlisting}

% =============================================================================
\section{Discussion and Expected Behaviour}
% =============================================================================

The on-board demonstration was reviewed live. For completeness, the
expected UART trace under three representative input patterns is:

\begin{lstlisting}[language={},basicstyle=\ttfamily\footnotesize,
                   numbers=none,backgroundcolor=\color{codebg},
                   frame=single,breaklines=true,
                   xleftmargin=6pt,xrightmargin=2pt]
# Happy path:  BTN0, BTN0, BTN1
[BTN0] S0 -> S1  LED=OFF
[BTN0] S1 -> S2  LED=OFF
[BTN1] S2 -> S3  LED=ON   *** DETECTED #1 ***
[BTN0] S3 -> S1  LED=OFF       (LED clears the moment we leave S3)

# KMP path:  BTN0, BTN0, BTN0, BTN1   (extra BTN0 keeps us at S2)
[BTN0] S0 -> S1  LED=OFF
[BTN0] S1 -> S2  LED=OFF
[BTN0] S2 -> S2  LED=OFF       (KMP self-loop, prefix preserved)
[BTN1] S2 -> S3  LED=ON   *** DETECTED #2 ***

# False start:  BTN0, BTN1
[BTN0] S0 -> S1  LED=OFF
[BTN1] S1 -> S0  LED=OFF
\end{lstlisting}

Three behaviours are worth highlighting. \textbf{(1) Moore output:}
the LED is a pure function of \texttt{state}, so it lights up the
instant the FSM enters \texttt{S3} and clears automatically on any
subsequent press. \textbf{(2) KMP awareness at \texttt{S2} and
\texttt{S3}:} an extra \texttt{BTN0} after \texttt{S2} or right after a
detection does not throw away the matched prefix --- it preserves the
\texttt{BTN0}-as-new-start-of-pattern semantics, which is what makes
back-to-back matches like \texttt{BTN0,BTN0,BTN1,BTN0,BTN0,BTN1}
produce two detections. \textbf{(3) Edge detection in software:}
the 5\,ms polling loop only reacts to a 0$\rightarrow$1 transition on
each button line, so holding a button does not flood the FSM with
spurious events; samples where both buttons are high simultaneously
are dropped to avoid ambiguous transitions.

% =============================================================================
\section{Conclusion}
% =============================================================================

The PS+PL system meets both deliverables: a bitstream \texttt{.xsa}
exported from Vivado pairs the Zynq PS with two AXI GPIO blocks for
buttons and LEDs, and the bare-metal application implements the
4-state KMP detector for \texttt{BTN0, BTN0, BTN1} with
\texttt{LD0\_G} as the detection indicator. The on-board demonstration
matched the expected trace above: \texttt{LD0\_G} lights up exactly
when the third press completes a valid sequence, redundant
\texttt{BTN0}s do not invalidate the in-progress match, and any other
press correctly resets the FSM.

\end{document}
